\documentclass[11pt]{article}
\usepackage[margin=1in]{geometry}
\usepackage{amsmath,amssymb,amsthm,mathtools}
\usepackage{graphicx}
\usepackage{hyperref}
\usepackage{xcolor}

\newtheorem{theorem}{Theorem}
\newtheorem{lemma}[theorem]{Lemma}
\newtheorem{remark}[theorem]{Remark}
\newcommand{\R}{\mathbb R}
\newcommand{\Bc}{B^{b}_{c}}
\newcommand{\supp}{\operatorname{supp}}

\title{Equicontinuous Convergence Need Not Have the Lebesgue Property}
\author{Automatic mathematics run \texttt{fa\_banach\_001}}
\date{26 August 2026}

\begin{document}
\setlength{\emergencystretch}{3em}
\maketitle

\begin{center}
\fbox{\parbox{0.91\textwidth}{
\textbf{Status.} Candidate counterexample, likely valid, for human review.
It gives a negative answer to the general continuous-operator question on
printed page 8 of Zabeti's paper.  The domain is locally convex and has the
$AM$ and Levi properties; the range is order complete and Lebesgue.  The
failure is caused only by the equicontinuous convergence topology.
}}
\end{center}

\section{Source question}

The source is O.~Zabeti, \emph{$AM$-spaces from a locally solid vector lattice
point of view with applications}, arXiv:1912.00141, later published in the
\emph{Bulletin of the Iranian Mathematical Society} \cite{source}.
Theorem 5 of that paper proves that if $X$ has the $AM$ and Levi properties,
$Y$ is order complete, and $Y$ has the Lebesgue property, then the lattice of
order-bounded operators $B^b(X,Y)$ has the Lebesgue property for the topology
of uniform convergence on bounded sets.  On printed page 8 the author asks
whether a version holds for continuous order-bounded operators with the
equicontinuous convergence topology.  The paper proves the positive result
when $X$ is Frechet (Corollary 7).

\begin{center}
\includegraphics[width=0.97\textwidth]{figures/open_problem_crop.png}
\end{center}

Thus the natural general question is:

\begin{quote}
If $X$ and $Y$ are locally solid vector lattices, $X$ has the $AM$ and Levi
properties, $Y$ is order complete and has the Lebesgue property, must
$\Bc(X,Y)$ have the Lebesgue property for equicontinuous convergence?
\end{quote}

Here a net $(S_\alpha)$ converges equicontinuously to zero if, for every zero
neighborhood $V$ in $Y$, there is a zero neighborhood $U$ in $X$ such that,
for every $\varepsilon>0$, eventually
\[
   S_\alpha(U)\subseteq \varepsilon V.
\]

\section{Idea of the proof}

Use an uncountable coordinate set $I$.  Give the finite-support lattice
$\R^{(I)}$ every weighted $\ell^1$ seminorm.  This makes bounded sets extremely
small: each bounded set uses only finitely many coordinates.  The domain
therefore has both the $AM$ and Levi properties.

For each finite $F\subset I$, delete the coordinates in $F$ and map the
remaining finite-support vector into $\ell^1(I)$.  These positive continuous
operators decrease to zero, because each individual input has finite support.
Equicontinuous convergence would require one domain neighborhood on which all
sufficiently late coordinate deletions are uniformly small.  A neighborhood
uses only finitely many weight functions.  On an uncountable set, their
combined reciprocal cannot tend to zero outside finite sets.  Hence every
finite deletion leaves a coordinate vector of one fixed positive size inside
the neighborhood, and its image has the same fixed $\ell^1$ norm.

\section{The counterexample}

Let $I$ be an uncountable set and put
\[
   X=\R^{(I)}
   =\{x\in\R^I:\supp(x)\text{ is finite}\},
\]
with coordinatewise order.  For every weight $a=(a_i)_{i\in I}\in[0,\infty)^I$
define
\[
   p_a(x)=\sum_{i\in I}a_i|x_i|.
\]
The sum is finite for every $x\in X$.  Let $\tau$ be the locally convex
topology generated by all seminorms $p_a$.  Since each $p_a$ is a lattice
seminorm and the coordinate weights separate points, $(X,\tau)$ is a
Hausdorff locally convex locally solid vector lattice.

\begin{lemma}\label{lem:bounded}
A subset $B\subset X$ is $\tau$-bounded if and only if there is a finite
$J\subset I$ such that $B\subset\R^J$ and $B$ is bounded in that
finite-dimensional space.
\end{lemma}

\begin{proof}
The reverse implication is immediate.  Conversely, suppose that the union of
the supports of elements of $B$ is infinite.  Choose $x_n\in B$ and pairwise
distinct $i_n\in I$ with $x_n(i_n)\ne0$.  Define one weight by
\[
   a_{i_n}=\frac{n}{|x_n(i_n)|},\qquad a_i=0
   \quad(i\notin\{i_n:n\ge1\}).
\]
Then $p_a(x_n)\ge n$, contradicting boundedness.  Thus all supports lie in a
finite $J$.  The coordinate seminorms show that every coordinate of $B$ is
bounded, which is equivalent to finite-dimensional boundedness.
\end{proof}

\begin{lemma}\label{lem:domain}
The locally solid vector lattice $X$ has the $AM$ and Levi properties.
\end{lemma}

\begin{proof}
Let $B$ be bounded.  By Lemma~\ref{lem:bounded}, all its elements have support
in one finite set $J$ and are coordinatewise bounded.  Every finite lattice
supremum of elements of $B$ has the same two properties, so $B^\vee$ is
bounded.  This is the $AM$ property.

If $D\subset X_+$ is bounded and upward directed, it is contained in
$\R^J$ for some finite $J$.  Its coordinatewise supremum is finite in each
coordinate and belongs to $\R^J\subset X$.  It is the supremum of $D$, which
proves the Levi property.
\end{proof}

Let $Y=\ell^1(I)$ with its norm topology and coordinatewise order.

\begin{lemma}\label{lem:range}
The Banach lattice $Y$ is order complete and has the Lebesgue property.
\end{lemma}

\begin{proof}
If a nonempty subset of $Y$ is bounded above by $u\in Y$, its coordinatewise
supremum differs from any fixed member by a positive vector dominated by an
$\ell^1$ vector.  Hence the supremum belongs to $Y$.  Thus $Y$ is order
complete.

Now let $y_\beta\downarrow0$ in $Y$.  Fix $\beta_0$.  The support of
$y_{\beta_0}$ is countable, and $0\le y_\beta\le y_{\beta_0}$ for
$\beta\ge\beta_0$.  Given $\eta>0$, choose a finite $K\subset I$ for which
\[
   \sum_{i\notin K}y_{\beta_0}(i)<\frac{\eta}{2}.
\]
Coordinatewise decrease to zero and directedness give $\beta_1\ge\beta_0$
such that
\[
   \sum_{i\in K}y_\beta(i)<\frac{\eta}{2}
   \quad(\beta\ge\beta_1).
\]
Therefore $\|y_\beta\|_1<\eta$ eventually.  This is the Lebesgue property.
\end{proof}

Direct the finite subsets $F\subset I$ by inclusion and define
\[
   T_F:X\longrightarrow Y,
   \qquad (T_Fx)_i=
   \begin{cases}
      0,&i\in F,\\
      x_i,&i\notin F.
   \end{cases}
\]

\begin{theorem}\label{thm:main}
Every $T_F$ is positive, continuous, and order bounded, and
\[
   T_F\downarrow0\quad\text{in }\Bc(X,Y).
\]
However, $(T_F)$ does not converge to zero in the equicontinuous convergence
topology.  Consequently $\Bc(X,Y)$ does not have the Lebesgue property.
\end{theorem}

\begin{proof}
Positivity is clear, and every positive linear operator is order bounded.  If
$\mathbf 1_{I\setminus F}$ is the $0$--$1$ weight of the complement of $F$,
then
\[
   \|T_Fx\|_1=p_{\mathbf 1_{I\setminus F}}(x),
\]
so $T_F$ is continuous.  If $F\subset G$, then $0\le T_G\le T_F$.  If a
positive operator $S$ is a lower bound for every $T_F$, then, for each
$i\in I$, choosing $F=\{i\}$ gives
\[
   0\le S e_i\le T_{\{i\}}e_i=0.
\]
Since the coordinate vectors span $X$, $S=0$.  Hence $T_F\downarrow0$ in the
operator lattice.

It remains to disprove equicontinuous convergence.  Let
\[
   V=\{y\in Y:\|y\|_1<1\}.
\]
Take an arbitrary zero neighborhood $U$ in $X$.  It contains a basic
neighborhood
\[
   W=\{x:p_{a^j}(x)<\delta_j\text{ for }j=1,\ldots,m\},
\]
where the $a^j$ are nonnegative weights and $\delta_j>0$.  For every $i\in I$
put
\[
   b_i=\max\left\{1,
      \frac{a^1_i}{\delta_1},\ldots,
      \frac{a^m_i}{\delta_m}\right\}.
\]
Every $b_i$ is finite and
\[
   I=\bigcup_{M=1}^{\infty}\{i\in I:b_i\le M\}.
\]
Because $I$ is uncountable, for some integer $M$ the set
$A_M=\{i:b_i\le M\}$ is infinite.  For every finite $F\subset I$, choose
$i\in A_M\setminus F$ and set $x=(2M)^{-1}e_i$.  Then, for each $j$,
\[
   p_{a^j}(x)=\frac{a^j_i}{2M}\le\frac{\delta_j}{2},
\]
so $x\in W\subset U$.  But $i\notin F$, and therefore
\[
   \|T_Fx\|_1=\frac{1}{2M}.
\]
With $\varepsilon=(4M)^{-1}$, this shows
$T_F(U)\not\subset\varepsilon V$ for every finite $F$.  Thus there is no
zero neighborhood $U$ satisfying the defining eventual-containment condition
for equicontinuous convergence.
\end{proof}

\begin{remark}
The uncountable direct-sum topology is not metrizable, so the counterexample
does not conflict with the source paper's positive Frechet result.  The proof
also identifies the precise obstruction: one fixed neighborhood cannot encode
decay along uncountably many coordinates modulo finite sets.
\end{remark}

\section{Verification and limitations}

The construction is purely exact and uses no numerical computation.  The
hypotheses were audited separately: Lemma~\ref{lem:bounded} proves the needed
bounded-set structure, Lemma~\ref{lem:domain} verifies both domain properties,
Lemma~\ref{lem:range} treats arbitrary decreasing nets in the range, and
Theorem~\ref{thm:main} checks both order convergence and the literal
quantifiers in equicontinuous convergence.

The counterexample answers the unrestricted question negatively.  It does not
weaken Corollary 7 of the source: additional metrizability or Frechet
hypotheses still give the positive conclusion established there.  Human review
should focus on whether the phrase ``a version of Theorem 5'' is intended with
the printed hypotheses only; under that natural reading, every hypothesis is
met.

\section{Bounded novelty check}

On 26 August 2026, the four lightweight indexes of run
\texttt{fa\_banach\_001} were searched for the arXiv id, title, and operator
notation.  Exact-phrase, title, and close-variant arXiv/web searches were also
performed.  Two related later arXiv papers on locally solid operator
topologies and Lebesgue operators were checked at title/abstract level
\cite{dengdejeu,alpay}; neither was found to state or resolve this exact
question.  No later answer was found within these bounded searches.  This is
evidence of plausible novelty, not a certification.

\enlargethispage{2\baselineskip}
\footnotesize
\begin{thebibliography}{9}

\bibitem{source}
O.~Zabeti,
\emph{$AM$-spaces from a locally solid vector lattice point of view with
applications}, arXiv:1912.00141; Bull. Iran. Math. Soc. 47 (2021),
\href{https://doi.org/10.1007/s41980-020-00458-7}{doi:10.1007/s41980-020-00458-7}.

\bibitem{dengdejeu}
Y.~Deng and M.~de Jeu,
\emph{Convergence structures and locally solid topologies on vector lattices
of operators}, arXiv:2008.05379.

\bibitem{alpay}
S.~Alpay, E.~Emelyanov, and S.~Gorokhova,
\emph{The $o\tau$-continuous, Lebesgue, KB, and Levi operators between vector
lattices and topological vector spaces}, arXiv:2105.01810.

\end{thebibliography}

\end{document}
